\NativeTitle{29}{第四单元素养评价（B）}{（运算律）}{45}{5}{一 & 二 & 三 & 四 & 五}{ &  &  &  & }
\NativeSlot{282}{56}{%
\ExamGroup{一、填空。{\normalfont\kaishu\fontsize{10.8bp}{14bp}\selectfont （每空1分，共27分）}}

% p029_s01

\ExamQuestion{1.}{5}%
超市运来3箱文具，第一箱有277件，第二箱有164件，第三箱有536件。计算一共有多少件文具时，先把第\mbox{（\hspace{1.8em}）}箱和第\mbox{（\hspace{1.8em}）}箱的数量合起来，再加上第\mbox{（\hspace{1.8em}）}箱的数量，这样计算比较简便，列式为\mbox{（\hspace{1.8em}）}。\par
}
\NativeSlot{381}{36}{%
% p029_s02

\ExamQuestion{2.}{5}%
计算\(125\times5\times8\times20\)时，可以先算\mbox{（\hspace{1.8em}）}和\mbox{（\hspace{1.8em}）}（填算式），再算\mbox{（\hspace{1.8em}）}（填算式），这是运用了\mbox{（\hspace{1.8em}）}律和\mbox{（\hspace{1.8em}）}律使计算简便。\par
}
\NativeSlot{433}{65}{%
\ExamQuestion{3.}{5}%
在圆圈里填上“＞”“＜”或“＝”。\par\smallskip
\noindent%
% p029_s03
\makebox[.499\linewidth][l]{\(260\times48\ \bigcirc\ 26\times480\)}%
% p029_s04
\makebox[.499\linewidth][l]{\(301\times24\ \bigcirc\ 300\times24+1\)}\par
\noindent%
% p029_s05
\makebox[.499\linewidth][l]{\(145-18-82\ \bigcirc\ 145-(18+82)\)}%
% p029_s06
\makebox[.499\linewidth][l]{\(50\times(33-13)\ \bigcirc\ 50\times33-13\)}\par
\noindent%
% p029_s07
\makebox[.499\linewidth][l]{\((13+4)\times6\ \bigcirc\ 13\times6+4\times6\)}%
% p029_s08
\makebox[.499\linewidth][l]{\(50\div(2\times5)\ \bigcirc\ 50\div2\times5\)}\par
}
\NativeSlot{547}{36}{%
% p029_s09

\ExamQuestion{4.}{5}%
实验小学运动会开幕式有8个方队，每个方队25人，每人拿2个气球，一共需要\mbox{（\hspace{1.8em}）}个气球。\par
}
\NativeSlot{597}{58}{%
% p029_s10

\ExamQuestion{5.}{5}%
一块菜地每行能种18棵白菜，一共有98行，这块菜地共能种多少棵白菜？小明列式\(18\times98\)，结合右图，还可列综合算式\mbox{（\hspace{1.8em}）}，这是运用\mbox{（\hspace{1.8em}）}律。\par
\BookFigure{figures/p029_s10-01.tex}
}
\NativeSlot{696}{58}{%
% p029_s11

\ExamQuestion{6.}{5}%
明明用乘法分配律比较\(181\times67\)和\(182\times66\)的大小。\par\noindent （1）请帮他把思考和比较过程填完整。\(181\times67=181\times(\square+1)=181\times66+\square\times1\)。\par
}
